\documentclass[12pt,a4paper]{article}

% =========================================================
% Pacotes
% =========================================================
\usepackage[utf8]{inputenc}
\usepackage[T1]{fontenc}
\usepackage[brazil]{babel}
\usepackage{geometry}
\usepackage{amsmath}
\usepackage{amssymb}
\usepackage{booktabs}
\usepackage{array}
\usepackage{longtable}
\usepackage{xcolor}
\usepackage{listings}
\usepackage{enumitem}
\usepackage{hyperref}
\usepackage{float}
\usepackage{caption}

\geometry{
    left=2.5cm,
    right=2.5cm,
    top=2.5cm,
    bottom=2.5cm
}

\hypersetup{
    colorlinks=true,
    linkcolor=blue,
    citecolor=blue,
    urlcolor=blue
}

% =========================================================
% Configuração para códigos-fonte
% =========================================================
\definecolor{codegray}{rgb}{0.96,0.96,0.96}
\definecolor{codeframe}{rgb}{0.8,0.8,0.8}
\definecolor{codegreen}{rgb}{0.13,0.54,0.13}
\definecolor{codeblue}{rgb}{0.0,0.2,0.7}
\definecolor{codepurple}{rgb}{0.5,0.0,0.5}

\lstset{
    basicstyle=\ttfamily\footnotesize,
    backgroundcolor=\color{codegray},
    rulecolor=\color{codeframe},
    frame=single,
    breaklines=true,
    showstringspaces=false,
    tabsize=4,
    numbers=left,
    numberstyle=\tiny\color{gray},
    numbersep=7pt,
    columns=fullflexible,
    keepspaces=true
}

% =========================================================
% Início do Documento
% =========================================================
\begin{document}

\begin{titlepage}
    \centering

    {\Large Universidade Federal de Itajubá -- UNIFEI\par}
    \vspace{0.4cm}

    {\large Campus Avançado de Itabira\par}
    \vspace{0.4cm}

    {\large Engenharia de Computação\par}
    \vspace{1.2cm}

    {\Large ECOI2219 -- Compiladores\par}
    \vspace{1.5cm}

    {\Huge \textbf{Relatório Técnico -- Laboratório 1}\par}
    \vspace{0.5cm}

    {\LARGE Tradução Dirigida por Sintaxe com PLY\par}
    \vspace{0.4cm}
    {\large Da Análise Léxica à Máquina Virtual MEPA 5.0\par}

    \vfill

    \begin{flushleft}
        \textbf{Integrantes da Equipe:}\\
        Aluno 1\\
        Aluno 2\\[0.5cm]

        \textbf{Docente Responsável:} Prof. Walter Aoiama Nagai\\
        \textbf{Disciplina:} ECOI2219 -- Compiladores
    \end{flushleft}

    \vfill

    Itabira -- MG\\
    Outubro de 2026

\end{titlepage}

\tableofcontents
\newpage

% =========================================================
\section{Introdução}
% =========================================================

A compilação de linguagens de programação fundamenta-se na decomposição do
processo de tradução em etapas rigorosamente estruturadas: análise léxica,
análise sintática, análise semântica e geração de código intermediário ou de
máquina. Enquanto abordagens tradicionais frequentemente constroem uma Árvore
Sintática Abstrata (AST) intermediária completa antes de percorrer seus nós
para emitir código, a técnica de \textbf{Tradução Dirigida por Sintaxe}
(\textit{Syntax-Directed Translation} -- SDT) propõe um acoplamento direto:
as ações semânticas de verificação e emissão de código são executadas
concomitantemente com as reduções da gramática livre de contexto.

Este trabalho documenta o projeto, a implementação, a validação experimental e
a auditoria técnica do compilador para a linguagem didática $\mu$C, orientado à
emissão de instruções para a Máquina Virtual à Pilha (MEPA 5.0). O sistema foi
reconstruído integralmente em Python com o auxílio da biblioteca PLY
(\textit{Python Lex-Yacc}), cobrindo:

\begin{enumerate}[label=\textbf{\arabic*.}]
    \item \textbf{Passo 1 -- Analisador Léxico:} Reconhecimento de identificadores,
          números inteiros, números reais em ponto flutuante e notação científica,
          operadores relacionais/lógicos, comentários de linha e bloco, e
          mapeamento de palavras reservadas;
    \item \textbf{Passo 2 -- Tradução de Expressões:} Emissão em pós-ordem de
          operações aritméticas sobre a máquina de pilha, respeitando precedência
          e associatividade por meio de gramática estratificada em níveis;
    \item \textbf{Passo 3 -- Declarações, Memória e Atribuição:} Gerenciamento de
          tabela de símbolos, cálculo de deslocamentos (\textit{offsets}) em bytes
          para os tipos primitivos \texttt{int} (4 bytes) e \texttt{float} (8 bytes),
          atribuição de variáveis e escrita (\texttt{print});
    \item \textbf{Passo 4 -- Comandos Condicionais:} Resolução do problema do
          ``rótulo no meio'' (\textit{mid-rule actions}) por meio de produções
          auxiliares e propagação de atributos herdados com pilhas de rótulos;
    \item \textbf{Passo 5 -- Laços de Repetição e Desvios:} Padrão de repetição
          \texttt{while} e implementação do comando \texttt{break} sensível ao
          escopo do laço mais interno;
    \item \textbf{Desafios 4.1 a 4.4:} Extensão com comando pós-testado
          \texttt{do-while}, suporte a leitura com \texttt{read} e blocos aninhados
          com preservação de alocação no prólogo, mapeamento estruturado com a
          implementação Java de referência, e auditoria formal de código gerado por IA.
\end{enumerate}

Todas as decisões arquiteturais foram confrontadas empiricamente contra o
compilador de referência em Java da disciplina (\texttt{uc-compilador.jar}) e
validadas por execução direta no interpretador oficial MEPA 5.0.

% =========================================================
\section{Ambiente de Desenvolvimento}
% =========================================================

O ambiente de desenvolvimento e execução adotado pela equipe reúne:

\begin{itemize}[leftmargin=1.5cm]
    \item \textbf{Linguagem Base:} Python 3.13 (compatível com especificações Python $\ge$ 3.10);
    \item \textbf{Gerador Léxico/Sintático:} Biblioteca PLY (\textit{Python Lex-Yacc}) versão 3.11;
    \item \textbf{Máquina Virtual Alvo:} Interpretador MEPA 5.0 (\texttt{mepa\_pt.py} / \texttt{mepa.py}),
          de autoria do Prof. Tomasz Kowaltowski;
    \item \textbf{Compilador de Referência:} Implementação Java em analisador sintático
          descendente recursivo (\texttt{uc-full/bin} e \texttt{uc-compilador.jar}),
          executada sob Java SE Runtime Environment 25 LTS;
    \item \textbf{Formatador e Tipografia:} Sistema \LaTeX{} (MiKTeX 25.4 / pdf\TeX{} 4.21).
\end{itemize}

A estrutura de arquivos do projeto está organizada como segue:

\begin{lstlisting}[numbers=none]
lab/
|-- lab1_tds.py          # Codigo-fonte completo do compilador
|-- modulo4-lab1.pdf     # Guia oficial de requisitos da atividade
|-- main.tex             # Codigo-fonte do relatorio tecnico em LaTeX
|-- exemplos/            # Suite de testes em linguagem uC (.uc)
|   |-- expr.uc          # Expressoes basicas e atribuicao
|   |-- if_else.uc       # Desvio condicional if-else
|   |-- while.uc         # Repeticao pre-testada
|   |-- break.uc         # Interrupcao controlada por break
|   |-- do_while.uc      # Desafio 4.1: repeticao pos-testada
|   |-- fib.uc           # Desafio 4.2: Fibonacci, read e blocos aninhados
|   |-- erro_var.uc      # Rejeicao semantica de variavel nao declarada
|   `-- erro_break.uc    # Rejeicao semantica de break fora de laco
|-- mepa/                # Interpretador MEPA 5.0 oficial
`-- uc-full/             # Compilador uC de referencia (Java / dist / bin)
\end{lstlisting}

% =========================================================
\section{Passo 1 -- Analisador Léxico}
% =========================================================

\subsection{Princípios e Regras Léxicas}

O analisador léxico (\textit{scanner}) tem a responsabilidade primária de transformar
a sequência contínua de caracteres do arquivo-fonte em uma sequência de unidades
significativas denominadas \textit{tokens}.

No PLY, o analisador léxico é instanciado através do módulo \texttt{ply.lex}. O
módulo emprega introspecção em Python: variáveis com o prefixo \texttt{t\_} que
contêm expressões regulares ou funções definem o padrão dos tokens.

\subsection{Palavras Reservadas e Identificadores}

Para evitar colisões léxicas e proliferação desnecessária de expressões regulares,
adotou-se o padrão clássico da tabela de palavras-chave: todos os identificadores
são capturados pela expressão regular de \texttt{t\_ID}, e em seguida o lexema é
consultado no dicionário \texttt{reserved}:

\begin{lstlisting}[language=Python]
reserved = {
    "if": "IF",
    "else": "ELSE",
    "while": "WHILE",
    "do": "DO",          # Desafio 4.1
    "read": "READ",      # Desafio 4.2
    "break": "BREAK",
    "print": "PRINT",
    "true": "TRUE",
    "false": "FALSE",
    "int": "INT",
    "float": "FLOAT",
}

def t_ID(t):
    r"[a-zA-Z_][a-zA-Z0-9_]*"
    t.type = reserved.get(t.value, "ID")
    return t
\end{lstlisting}

Essa decisão garante que identificadores válidos como \texttt{ifx} ou \texttt{while\_count}
permaneçam classificados como \texttt{ID}, enquanto o lexema exato \texttt{if} é promovido
ao token \texttt{IF}.

\subsection{Operadores e a Regra do Maximal Munch}

Os operadores relacionais, lógicos e de atribuição foram definidos conforme a tabela
oficial do roteiro:

\begin{center}
\begin{tabular}{cclc}
\toprule
\textbf{Token} & \textbf{Texto} & \textbf{Expressão Regular no PLY} & \textbf{Significado} \\
\midrule
\texttt{EQ}     & \texttt{==}  & \texttt{r"=="}                  & Igualdade relacional \\
\texttt{NE}     & \texttt{!=}  & \texttt{r"!="}                  & Diferença relacional \\
\texttt{LE}     & \texttt{<=}  & \texttt{r"<="}                  & Menor ou igual \\
\texttt{GE}     & \texttt{>=}  & \texttt{r">="}                  & Maior ou igual \\
\texttt{AND}    & \texttt{\&\&}& \texttt{r"\&\&"}                & Conjunção lógica \\
\texttt{OR}     & \texttt{||}  & \texttt{r"\textbackslash|\textbackslash|"} & Disjunção lógica \\
\texttt{NOT}    & \texttt{!}   & \texttt{r"!"}                   & Negação unária \\
\texttt{LT}     & \texttt{<}   & \texttt{r"<"}                   & Menor que \\
\texttt{GT}     & \texttt{>}   & \texttt{r">"}                   & Maior que \\
\texttt{ASSIGN} & \texttt{=}   & \texttt{r"="}                   & Atribuição \\
\bottomrule
\end{tabular}
\end{center}

\noindent\textbf{O princípio do Maximal Munch:} Quando o compilador processa uma cadeia
como \texttt{<=} ou \texttt{==}, o PLY classifica internamente as expressões regulares
de comprimento fixo por ordem decrescente de tamanho da expressão. Dessa forma,
a expressão \texttt{r"<="} (2 caracteres) tem prioridade determinística sobre
\texttt{r"<"} (1 caractere). O lexer consome a maior cadeia de caracteres possível que
forme um token válido, impedindo que \texttt{<=} seja fragmentado erroneamente em
\texttt{<} seguido de \texttt{=}.

\subsection{Tratamento Numérico: Inteiros e Notação Científica}

O roteiro prevê o reconhecimento de inteiros (\texttt{NUM}) e números reais (\texttt{REAL}):

\begin{lstlisting}[language=Python]
def t_REAL(t):
    r"\d+\.\d+([eE][+-]?\d+)?|\d+[eE][+-]?\d+"
    t.value = float(t.value)
    return t

def t_NUM(t):
    r"\d+"
    t.value = int(t.value)
    return t
\end{lstlisting}

A expressão regular adotada para \texttt{t\_REAL} foi expandida para suportar
tanto a representação fracionária padrão (\texttt{3.14}) quanto notação exponencial
comum em programas científicos (\texttt{3.14e-2}, \texttt{1e5}), alinhando-se
perfeitamente com o compilador Java de referência (\texttt{Lexer.java}).

\subsection{Comentários, Linhas e Erros Léxicos}

Espaços em branco, quebras de linha e comentários devem ser descartados sem
gerar tokens sintáticos, mas mantendo a numeração precisa das linhas:

\begin{lstlisting}[language=Python]
t_ignore = " \t\r"

def t_newline(t):
    r"\n+"
    t.lexer.lineno += len(t.value)

def t_COMMENT_LINE(t):
    r"//[^\n]*"
    pass

def t_COMMENT_BLOCK(t):
    r"/\*[\s\S]*?\*/"
    t.lexer.lineno += t.value.count("\n")

def t_error(t):
    print(f"erro lexico na linha {t.lexer.lineno}: caractere ilegal {t.value[0]!r}", file=sys.stderr)
    t.lexer.skip(1)
\end{lstlisting}

\subsection{Validação Experimental do Analisador Léxico}

O teste proposto no roteiro foi executado com o comando:
\begin{lstlisting}[language=bash]
python lab1_tds.py exemplos/expr.uc --tokens
\end{lstlisting}

\noindent O programa de teste \texttt{exemplos/expr.uc} contém:
\begin{lstlisting}[language=C]
{
    int a;
    a = 1 + 2;
}
\end{lstlisting}

\noindent A saída obtida no console correspondeu exatamente à especificação:
\begin{lstlisting}
  1 {          '{'
  2 INT        'int'
  2 ID         'a'
  2 ;          ';'
  3 ID         'a'
  3 ASSIGN     '='
  3 NUM        1
  3 +          '+'
  3 NUM        2
  3 ;          ';'
  4 }          '}'
\end{lstlisting}

Observa-se que:
\begin{itemize}
    \item \texttt{int} foi reconhecido como \texttt{INT} (palavra reservada), e não como \texttt{ID};
    \item Os identificadores \texttt{a} foram reconhecidos como \texttt{ID};
    \item Os inteiros \texttt{1} e \texttt{2} foram convertidos para inteiros nativos do Python;
    \item As contagens de linhas acompanharam perfeitamente o arquivo-fonte.
\end{itemize}

% =========================================================
\section{Passo 2 -- Tradução de Expressões}
% =========================================================

\subsection{Teoria da Tradução Dirigida por Sintaxe (SDT)}

Na tradução dirigida por sintaxe associada a analisadores do tipo \textit{bottom-up}
(como o LALR(1) gerado pelo PLY/Yacc), as ações semânticas associadas a cada regra
são disparadas exclusivamente quando a produção atinge sua \textbf{redução}.

Considere a gramática tradicional com atributos sintetizados:
\begin{align*}
    E &\rightarrow E_1 + T \quad & \{ E.val = E_1.val + T.val \} \\
    E &\rightarrow T           & \{ E.val = T.val \} \\
    T &\rightarrow \text{\textbf{num}} & \{ T.val = \text{\textbf{num}}.lexval \}
\end{align*}

Num compilador para máquina de pilha (MEPA), o objetivo não é calcular o valor numérico
em tempo de compilação, mas emitir as instruções que colocarão os operandos na pilha
e em seguida dispararão a operação correspondente.

\subsection{Por que a Ordem Funciona: Pós-Ordem na Árvore de Derivação}

Na análise ascendente LALR, as folhas da árvore sintática são reconhecidas e reduzidas
primeiro. Os nós intermediários reduzem depois dos nós filhos, e a raiz reduz por
último. Trata-se de uma travessia estrita em \textbf{pós-ordem}:
\begin{center}
\textbf{Filho Esquerdo} $\rightarrow$ \textbf{Filho Direito} $\rightarrow$ \textbf{Nó Pai}
\end{center}

Para a expressão:
\[
a + 2
\]
O autômato executa as seguintes reduções e emissões:
\begin{center}
\begin{tabular}{lll}
\toprule
\textbf{Redução} & \textbf{Instrução Emitida} & \textbf{Estado da Pilha MEPA} \\
\midrule
\texttt{factor : ID(a)}     & \texttt{CRVL 0,0} & $[a]$ \\
\texttt{factor : NUM(2)}    & \texttt{CRCT 2}   & $[a, 2]$ \\
\texttt{expr : expr + term} & \texttt{SOMA}     & $[a + 2]$ \\
\bottomrule
\end{tabular}
\end{center}

Como o operador \texttt{SOMA} desempilha os dois valores do topo da pilha, realiza a adição
e empilha o resultado, a pós-ordem natural das reduções do parser ascendente é
\textbf{rigorosamente a ordem esperada por uma máquina de pilha}.

\subsection{Resolução de Precedência por Níveis de Produção}

Para eliminar ambiguidades sintáticas e garantir associatividade à esquerda sem
necessidade de declarar prioridades artificiais, a gramática de expressões do $\mu$C
foi estruturada em quatro níveis hierárquicos:

\begin{align*}
    \text{expr}  &\rightarrow \text{expr } '+' \text{ term} \mid \text{expr } '-' \text{ term} \mid \text{term} \\
    \text{term}  &\rightarrow \text{term } '*' \text{ unary} \mid \text{term } '/' \text{ unary} \mid \text{unary} \\
    \text{unary} &\rightarrow '-' \text{ unary} \mid '+' \text{ unary} \mid \text{factor} \\
    \text{factor}&\rightarrow \text{\textbf{id}} \mid \text{\textbf{num}} \mid \text{\textbf{real}} \mid \text{\textbf{true}} \mid \text{\textbf{false}} \mid '(' \text{bool} ')'
\end{align*}

\noindent\textbf{Armadilha evitada:} A produção simplória \texttt{expr : expr '+' expr}
introduz ambiguidades e gera conflitos \textit{shift/reduce}. Por padrão, o Yacc/PLY
desempataria realizando o \textit{shift}, o que acarretaria associatividade à direita
(incorreta para subtração e divisão). A gramática em níveis força o agrupamento à esquerda:
para reduzir \texttt{a - b - c}, o compilador é obrigado a reduzir \texttt{a - b} primeiro.

\subsection{Implementação no PLY}

As funções correspondentes no arquivo \texttt{lab1\_tds.py} são implementadas de
forma modular:

\begin{lstlisting}[language=Python]
def p_expr_plus(p):
    "expr : expr '+' term"
    emit("SOMA")

def p_expr_minus(p):
    "expr : expr '-' term"
    emit("SUBT")

def p_term_times(p):
    "term : term '*' unary"
    emit("MULT")

def p_term_divide(p):
    "term : term '/' unary"
    emit("DIVI")

def p_unary_minus(p):
    "unary : '-' unary"
    emit("INVR")

def p_factor_num(p):
    "factor : NUM"
    emit(f"CRCT {p[1]}")

def p_factor_real(p):
    "factor : REAL"
    emit(f"CRCT {p[1]}")
\end{lstlisting}

\subsection{Verificação Experimental de Precedência}

Comparou-se a tradução de duas expressões com precedências distintas:

\begin{enumerate}
    \item \textbf{Sem parênteses:} \texttt{a + b * 3}
          \begin{lstlisting}
CRVL 0,0    ; a
CRVL 0,4    ; b
CRCT 3
MULT        ; multiplicacao ocorre primeiro
SOMA        ; soma ocorre por ultimo
          \end{lstlisting}
    \item \textbf{Com parênteses:} \texttt{(a + b) * 3}
          \begin{lstlisting}
CRVL 0,0    ; a
CRVL 0,4    ; b
SOMA        ; soma ocorre primeiro devido aos parenteses
CRCT 3
MULT        ; multiplicacao ocorre depois
          \end{lstlisting}
\end{enumerate}

O teste comprovou a aderência estrita à precedência algébrica padrão.

% =========================================================
\section{Passo 3 -- Declarações, Memória, Atribuição e Escrita}
% =========================================================

\subsection{Tabela de Símbolos e Deslocamento Global}

Para vincular identificadores a posições de memória, o compilador mantém o estado
de alocação através de:

\begin{lstlisting}[language=Python]
OFFSET = 0       # Proximo endereco disponivel na memoria
SYMBOLS = {}      # Tabela com o deslocamento de cada variavel
SCOPES = []       # Pilha de dicionarios para controle de escopo lexico
ALLOCATIONS = []  # Instrucoes AMEM coletadas para emissao no prologo
\end{lstlisting}

\subsection{Larguras de Tipos e Alocação de Memória}

Na linguagem $\mu$C, variáveis inteiras (\texttt{int}) possuem largura de 4 bytes,
enquanto números reais (\texttt{float}) possuem largura de 8 bytes (conforme a
classe \texttt{symbols/Type.java} do compilador de referência):

\begin{lstlisting}[language=Python]
def p_decl_int(p):
    "decl : INT ID ';'"
    declare_symbol(p[2], 4, p.lineno(2))

def p_decl_float(p):
    "decl : FLOAT ID ';'"
    declare_symbol(p[2], 8, p.lineno(2))
\end{lstlisting}

A função \texttt{declare\_symbol} insere o identificador no escopo ativo, armazena
o \textit{offset} atual e incrementa a contagem de memória:

\begin{lstlisting}[language=Python]
def declare_symbol(name, width, line):
    global OFFSET
    if name in SCOPES[-1]:
        semantic_error(f"linha {line}: variavel {name!r} ja declarada no bloco")
    SCOPES[-1][name] = OFFSET
    SYMBOLS[name] = OFFSET
    ALLOCATIONS.append(f"AMEM {width}\t; {name} address: {OFFSET}")
    OFFSET += width
\end{lstlisting}

\noindent\textbf{Princípio do Offset Sequencial Global:} O \textit{offset} jamais é
reiniciado na abertura de blocos aninhados \texttt{\{ \dots \}}. A numeração sequencial
é mantida para garantir que nenhuma variável sobrescreva a célula de memória de outra.
Para o trecho:
\begin{lstlisting}[language=C]
{
    int a;
    float x;
}
\end{lstlisting}
O compilador emite:
\begin{lstlisting}
AMEM 4    ; a address: 0
AMEM 8    ; x address: 4
\end{lstlisting}
Totalizando 12 bytes alocados e resultando em $\text{OFFSET} = 12$.

\subsection{Atribuição, Leitura de Variáveis e Escrita}

A atribuição requer que o valor resultante da expressão seja armazenado na célula
de memória correspondente:

\begin{lstlisting}[language=Python]
def p_stmt_assign(p):
    "stmt : ID ASSIGN bool ';'"
    address = lookup_symbol(p[1], p.lineno(1))
    emit(f"ARMZ 0,{address}\t; {p[1]}")

def p_factor_id(p):
    "factor : ID"
    address = lookup_symbol(p[1], p.lineno(1))
    emit(f"CRVL 0,{address}\t; {p[1]}")

def p_stmt_print(p):
    "stmt : PRINT '(' bool ')' ';'"
    emit("IMPR")
\end{lstlisting}

A busca de identificadores (\texttt{lookup\_symbol}) percorre a pilha \texttt{SCOPES}
da posição mais profunda para a mais externa:

\begin{lstlisting}[language=Python]
def lookup_symbol(name, line):
    for scope in reversed(SCOPES):
        if name in scope:
            return scope[name]
    semantic_error(f"linha {line}: variavel {name!r} nao declarada ou fora do escopo")
\end{lstlisting}

Se a variável não for encontrada, o compilador dispara um erro semântico amigável
com coordenadas de linha e encerra a compilação via \texttt{SystemExit(1)}, impedindo
a propagação de exceções brutas do Python como \texttt{KeyError}.

% =========================================================
\section{Passo 4 -- Comando Condicional e o Problema do Rótulo no Meio}
% =========================================================

\subsection{O Problema Identificado}

O padrão semântico esperado para uma instrução condicional \texttt{if} é:
\begin{lstlisting}[numbers=none]
    [codigo que avalia a condicao]
    DSVF L1        ; se falso, desvia para L1
    [codigo do corpo do comando]
L1: NADA
\end{lstlisting}

A tentativa ingênua de codificar essa estrutura no PLY consiste em escrever:
\begin{lstlisting}[language=Python]
# TENTATIVA INCORRETA:
def p_stmt_if(p):
    "stmt : IF '(' bool ')' stmt"
    l1 = new_label()
    emit(f"DSVF {l1}")  # EMISSAO TARDE DEMAIS!
    emit(f"{l1}: NADA")
\end{lstlisting}

\noindent\textbf{Por que a tentativa ingênua falha criticamente?} No modelo de
análise ascendente LALR, uma função sintática só é disparada após \textbf{todos}
os símbolos do corpo da produção terem sido completamente reduzidos. Isso significa
que todo o código do corpo (\texttt{stmt}) já foi emitido antes que o corpo da função
\texttt{p\_stmt\_if} seja invocado.
Como consequência, a instrução \texttt{DSVF L1} seria emitida \textbf{depois} do corpo
que ela deveria desviar! O programa compilaria sem erros visíveis, porém na execução
o desvio jamais pularia o bloco \texttt{then}, desvirtuando por completo a semântica
do programa.

\subsection{A Solução: Regra Auxiliar com Atributo Herdado}

Para resolver o problema, é necessário que uma redução intermediária ocorra
\textbf{imediatamente após o teste da condição e antes da entrada no corpo}.
A produção principal é desmembrada em duas:

\begin{lstlisting}[language=Python]
IF_PENDING = []

def p_if_head(p):
    "if_head : IF '(' bool ')'"
    false_label = new_label()
    end_label = new_label()
    emit(f"DSVF {false_label}")  # Emitido exatamente no ponto correto!
    IF_PENDING.append((false_label, end_label))

def p_stmt_if(p):
    "stmt : if_head stmt %prec IFX"
    false_label, _ = IF_PENDING.pop()
    mark_label(false_label)
\end{lstlisting}

\noindent\textbf{A mecânica de atributos herdados via pilhas do módulo:}
Em teoria de compiladores, quando uma regra filha precisa de um valor gerado
anteriormente por uma regra irmã ou pai, esse dado é denominado \textbf{atributo herdado}.
Como o PLY só propaga atributos sintetizados naturalmente através de \texttt{p[0]},
o estado herdado viaja através de pilhas de variáveis no escopo do módulo
(\texttt{IF\_PENDING}).

\subsection{Comando \texttt{if-else} e Ambiguidade do Dangling Else}

Para a estrutura completa \texttt{if-else}, o padrão exige desviar o ramo \texttt{then}
ao término de sua execução para não invadir o ramo \texttt{else}:

\begin{lstlisting}[language=Python]
def p_else_mark(p):
    "else_mark :"
    false_label, end_label = IF_PENDING[-1]
    emit(f"DSVS {end_label}")      # Pula o ramo else ao fim do then
    mark_label(false_label)        # Ponto de entrada do ramo else

def p_stmt_if_else(p):
    "stmt : if_head stmt ELSE else_mark stmt"
    _, end_label = IF_PENDING.pop()
    mark_label(end_label)          # Ponto de juncao final
\end{lstlisting}

\noindent\textbf{Por que \texttt{IF\_PENDING} deve ser uma pilha?} Em construções com
condicionais aninhados como:
\begin{lstlisting}[language=C]
if (a > 0)
    if (b > 0)
        print(1);
\end{lstlisting}
o \texttt{if} interno instancia seus próprios rótulos. Se fosse utilizada uma variável
simples, os rótulos do \texttt{if} externo seriam irremediavelmente sobrescritos.
Com a estrutura em pilha (LIFO), o \texttt{if} mais interno consome o topo da pilha,
restaurando perfeitamente o contexto do \texttt{if} externo.

A ambiguidade clássica do \textit{dangling else} (onde o compilador precisa decidir
a qual \texttt{if} pertence uma cláusula \texttt{else} isolada) foi tratada atribuindo
precedência sintática via:
\begin{lstlisting}[language=Python]
precedence = (
    ("nonassoc", "IFX"),
    ("nonassoc", "ELSE"),
)
\end{lstlisting}
associando deterministicamente o \texttt{else} ao \texttt{if} imediatamente precedente.

\subsection{Validação Experimental do Condicional}

Utilizou-se o programa \texttt{exemplos/if\_else.uc}:
\begin{lstlisting}[language=C]
{
    int a;
    a = 5;
    if (a > 0)
        a = 1;
    else
        a = 2;
    print(a);
}
\end{lstlisting}

O código MEPA gerado por \texttt{lab1\_tds.py} foi:
\begin{lstlisting}
INPP
AMEM 4    ; a address: 0
CRCT 5
ARMZ 0,0  ; a
CRVL 0,0  ; a
CRCT 0
CMMA
DSVF L1
CRCT 1
ARMZ 0,0  ; a
DSVS L2
L1: NADA
CRCT 2
ARMZ 0,0  ; a
L2: NADA
CRVL 0,0  ; a
IMPR
PARA
FIM
\end{lstlisting}

\noindent\textbf{Resultados da execução no MEPA:}
\begin{itemize}
    \item Para $a = 5$ (condição verdadeira): executou 15 instruções e imprimiu \textbf{1};
    \item Para $a = -5$ (condição falsa): executou 16 instruções e imprimiu \textbf{2}.
\end{itemize}
Ambos os ramos executaram com 100\% de fidelidade semântica.

% =========================================================
\section{Passo 5 -- Laços de Repetição e Desvios (\texttt{while} e \texttt{break})}
% =========================================================

\subsection{Padrão de Tradução do Laço \texttt{while}}

Diferentemente do \texttt{if}, em que o desvio aponta sempre para a frente, o laço
de repetição exige um salto para trás (\textit{back-jump}). Para que o comando possa
retornar à verificação da condição, o rótulo de início precisa ser emitido
\textbf{antes} da avaliação da expressão booleana.

O padrão implementado utiliza duas produções auxiliares:
\begin{lstlisting}[language=Python]
WHILE_PENDING = []
LOOP_EXITS = []

def p_while_start(p):
    "while_start :"
    start_label = new_label()
    exit_label = new_label()
    mark_label(start_label)        # Rotulo de retorno antes da condicao
    WHILE_PENDING.append((start_label, exit_label))
    LOOP_EXITS.append(exit_label)  # Registra saida para comandos break

def p_while_head(p):
    "while_head : WHILE '(' while_start bool ')'"
    _, exit_label = WHILE_PENDING[-1]
    emit(f"DSVF {exit_label}")     # Sai do laco se a condicao for falsa

def p_stmt_while(p):
    "stmt : while_head stmt"
    start_label, exit_label = WHILE_PENDING.pop()
    LOOP_EXITS.pop()
    emit(f"DSVS {start_label}")    # Volta para a condicao
    mark_label(exit_label)         # Rotulo de saida do laco
\end{lstlisting}

\subsection{Implementação do Comando \texttt{break}}

O comando \texttt{break} interrompe abruptamente o laço, transferindo o fluxo de
controle para o rótulo de término do laço mais interno. A pilha \texttt{LOOP\_EXITS}
armazena exatamente esses rótulos:

\begin{lstlisting}[language=Python]
def p_stmt_break(p):
    "stmt : BREAK ';'"
    if not LOOP_EXITS:
        semantic_error(f"linha {p.lineno(1)}: break fora de um laco")
    emit(f"DSVS {LOOP_EXITS[-1]}\t; break")
\end{lstlisting}

\noindent\textbf{Melhoria Documentada em Relação ao Compilador de Referência:}
Conforme alertado na Seção 10.2 do guia da disciplina, o compilador Java de referência
valida a ocorrência de \texttt{break} verificando a variável \texttt{loopDepth > 0},
porém \textbf{omite a geração da instrução DSVS correspondente} em \texttt{Parser.java}.
Como consequência, programas compilados pelo compilador de referência ignoravam o
\texttt{break} e executavam o laço até o fim.
O compilador desenvolvido pela nossa equipe corrige essa falha, emitindo explicitamente
\texttt{DSVS} para o rótulo de saída do laço corrente.

\subsection{Validação Experimental do Laço e do Break}

\subsubsection{Teste do Laço Padrão (\texttt{exemplos/while.uc})}

\begin{lstlisting}[language=C]
{
    int k;
    k = 1;
    while (k <= 5) {
        print(k);
        k = k + 1;
    }
}
\end{lstlisting}

O código MEPA gerado foi:
\begin{lstlisting}
INPP
AMEM 4    ; k address: 0
CRCT 1
ARMZ 0,0  ; k
L1: NADA  ; inicio do teste
CRVL 0,0  ; k
CRCT 5
CMEG
DSVF L2   ; sai para L2 se k > 5
CRVL 0,0  ; k
IMPR
CRVL 0,0  ; k
CRCT 1
SOMA
ARMZ 0,0  ; k
DSVS L1   ; reinicia laco
L2: NADA  ; saida
PARA
FIM
\end{lstlisting}

Executado no interpretador MEPA 5.0, realizou 71 instruções e produziu:
\begin{lstlisting}
1
2
3
4
5
\end{lstlisting}

\subsubsection{Teste de Interrupção por \texttt{break} (\texttt{exemplos/break.uc})}

\begin{lstlisting}[language=C]
{
    int k;
    k = 1;
    while (k <= 10) {
        if (k > 5)
            break;
        print(k);
        k = k + 1;
    }
}
\end{lstlisting}

O código MEPA gerado emite a instrução de salto do break:
\begin{lstlisting}
...
L1: NADA
CRVL 0,0    ; k
CRCT 10
CMEG
DSVF L2     ; saida do while
CRVL 0,0    ; k
CRCT 5
CMMA
DSVF L3     ; desvio do if
DSVS L2     ; break desvia diretamente para a saida L2!
L3: NADA
CRVL 0,0    ; k
IMPR
...
\end{lstlisting}

Executado no MEPA, o programa imprimiu exatamente os valores \texttt{1, 2, 3, 4, 5}
e encerrou a execução imediatamente após a ativação do desvio de interrupção.

\subsubsection{Teste de Rejeição de \texttt{break} Fora de Laço}

Para o arquivo \texttt{exemplos/erro\_break.uc}:
\begin{lstlisting}[language=C]
{
    int a;
    break;
}
\end{lstlisting}

A execução de \texttt{python lab1\_tds.py exemplos/erro\_break.uc} reportou:
\begin{lstlisting}
erro semantico: linha 3: break fora de um laco
\end{lstlisting}
comprovando a eficácia da guarda semântica baseada na pilha \texttt{LOOP\_EXITS}.

% =========================================================
\section{Execução no MEPA 5.0}
% =========================================================

\subsection{Arquitetura e Exigências Formais do MEPA}

O interpretador MEPA 5.0 (\texttt{mepa\_pt.py}), desenvolvido para o ensino de
compiladores, impõe exigências rígidas de formatação:
\begin{itemize}
    \item Todo programa deve iniciar com a instrução de carga \texttt{INPP} e
          terminar com as instruções \texttt{PARA} e \texttt{FIM};
    \item Instruções de manipulação de memória (\texttt{CRVL}, \texttt{ARMZ}) exigem
          dois argumentos separados por vírgula no formato $m,n$, onde $m$ representa
          o nível léxico e $n$ o deslocamento (como $\mu$C não possui sub-rotinas
          aninhadas, $m=0$ sempre);
    \item \textbf{Alinhamento de Rótulos:} O MEPA 5.0 exige obrigatoriamente que
          rótulos estejam na \textbf{mesma linha} de uma instrução executável. Linhas
          isoladas contendo apenas \texttt{L1:} provocam erro de sintaxe no carregador.
\end{itemize}

\subsection{Eliminação de Dependência Externa: O Padrão \texttt{mark\_label}}

No laboratório de referência em Java, o compilador emitia rótulos em linhas isoladas,
obrigando o uso de um script externo em \texttt{awk} ou \texttt{sed} para juntar
linhas antes de executar na MEPA:

\begin{lstlisting}[language=bash]
awk '/^[[:space:]]*L[0-9]+:[[:space:]]*$/ {p=$1; next}
{if (p!="") {sub(/^[[:space:]]*/,""); print p" "$0} else print; p=""}' mvpOutput.txt > mvp.mepa
\end{lstlisting}

A nossa implementação em Python superou essa limitação criando a função auxiliar:
\begin{lstlisting}[language=Python]
def mark_label(label):
    emit(f"{label}: NADA")
\end{lstlisting}

Como a instrução \texttt{NADA} (\textit{no-operation}) é uma instrução oficial
reconhecida pelo MEPA, os rótulos gerados pelo nosso compilador já saem prontos
para consumo imediato pelo interpretador, eliminando completamente a necessidade de
conversões intermediárias.

% =========================================================
\section{Desafios Implementados}
% =========================================================

\subsection{Desafio 4.1 -- Comando Pós-Testado \texttt{do-while}}

O comando pós-testado \texttt{do stmt while ( bool ) ;} difere estruturalmente do
\texttt{while} pré-testado porque o corpo é executado obrigatoriamente pelo menos
uma vez antes do primeiro teste lógico.

\subsubsection{Implementação Sintática e Semântica}

Por executar o corpo primeiro, não há o problema do ``rótulo no meio'': o rótulo
de início é posicionado antes do corpo, e a condição é testada no final:

\begin{lstlisting}[language=Python]
DO_PENDING = []

def p_do_head(p):
    "do_head : DO"
    start_label = new_label()
    exit_label = new_label()
    mark_label(start_label)
    DO_PENDING.append((start_label, exit_label))
    LOOP_EXITS.append(exit_label)

def p_stmt_do(p):
    "stmt : do_head stmt WHILE '(' bool ')' ';'"
    start_label, exit_label = DO_PENDING.pop()
    LOOP_EXITS.pop()
    emit(f"DSVF {exit_label}")   # Se falsa, encerra
    emit(f"DSVS {start_label}")   # Se verdadeira, reexecuta corpo
    mark_label(exit_label)
\end{lstlisting}

\subsubsection{Validação Experimental (\texttt{exemplos/do\_while.uc})}

\begin{lstlisting}[language=C]
{
    int i;
    i = 0;
    do {
        i = i + 1;
        print(i);
    } while (i < 3);
}
\end{lstlisting}

O compilador produziu a sequência MEPA:
\begin{lstlisting}
INPP
AMEM 4    ; i address: 0
CRCT 0
ARMZ 0,0  ; i
L1: NADA  ; corpo executado diretamente
CRVL 0,0  ; i
CRCT 1
SOMA
ARMZ 0,0  ; i
CRVL 0,0  ; i
IMPR
CRVL 0,0  ; i
CRCT 3
CMME      ; i < 3
DSVF L2
DSVS L1   ; volta ao inicio se i < 3
L2: NADA
PARA
FIM
\end{lstlisting}

Executado no interpretador MEPA 5.0, produziu a sequência:
\begin{lstlisting}
1
2
3
\end{lstlisting}
comprovando o pleno funcionamento da estrutura pós-testada.

\subsection{Desafio 4.2 -- Comando \texttt{read}, Blocos Aninhados e Alocação de Prólogo}

O Desafio 4.2 estabelece como critério de êxito compilar o algoritmo de Fibonacci
(\texttt{fib.uc}), alimentá-lo com a entrada $10$ no MEPA e obter como saída os
valores $10$ e $55$.

\subsubsection{Leitura de Dados e Escopos Aninhados}

O comando \texttt{read} lê um valor da entrada padrão e grava-o na memória:
\begin{lstlisting}[language=Python]
def p_stmt_read(p):
    "stmt : READ '(' ID ')' ';'"
    address = lookup_symbol(p[3], p.lineno(3))
    emit("LEIT")
    emit(f"ARMZ 0,{address}\t; {p[3]}")
\end{lstlisting}

Para suportar blocos aninhados com escopo léxico estático, as regras de bloco
gerenciam a pilha \texttt{SCOPES}:
\begin{lstlisting}[language=Python]
def p_block(p):
    "block : '{' enter_scope decls stmts '}'"
    leave_scope()

def p_enter_scope(p):
    "enter_scope :"
    enter_scope()

def p_stmt_block(p):
    "stmt : block"
    pass
\end{lstlisting}

\subsubsection{A Decisão Crítica de Alocação no Prólogo}

\noindent\textbf{O Problema da Reserva em Laço:} Se uma variável for declarada dentro
de um bloco interno subordinado a um \texttt{while} (como \texttt{\{ int temp; \dots \}}),
a emissão direta de \texttt{AMEM 4} durante a redução da declaração faria com que a
instrução \texttt{AMEM} fosse reexecutada a cada iteração do laço! Isso aumentaria
o ponteiro de topo da pilha indefinidamente, corrompendo as variáveis anteriores e
estourando a memória da máquina de pilha.

\noindent\textbf{A Solução Arquitetural:} O compilador separa a coleta de declarações
da emissão do código. Todas as instruções \texttt{AMEM} são registradas na lista
\texttt{ALLOCATIONS} e inseridas exclusivamente no \textbf{prólogo} do programa
(logo após \texttt{INPP}):
\begin{lstlisting}[language=Python]
def compile_source(source):
    ...
    parser.parse(source, lexer=lexer, tracking=True)
    return CODE[:1] + ALLOCATIONS + CODE[1:]
\end{lstlisting}
Dessa forma, o espaço total de memória necessário para o programa (incluindo variáveis
de blocos aninhados) é alocado uma única vez no início, garantindo a integridade dos
dados mesmo em repetições exaustivas.

\subsubsection{Validação Experimental com Fibonacci (\texttt{exemplos/fib.uc})}

\begin{lstlisting}[language=C]
{
    int n;
    int a;
    int b;
    int i;
    read(n);
    print(n);
    a = 0;
    b = 1;
    i = 0;
    while (i < n) {
        int temp;
        temp = a + b;
        a = b;
        b = temp;
        i = i + 1;
    }
    print(a);
}
\end{lstlisting}

\noindent O programa aloca todas as 5 variáveis no prólogo:
\begin{lstlisting}
INPP
AMEM 4    ; n address: 0
AMEM 4    ; a address: 4
AMEM 4    ; b address: 8
AMEM 4    ; i address: 12
AMEM 4    ; temp address: 16
LEIT
ARMZ 0,0  ; n
CRVL 0,0  ; n
IMPR
...
\end{lstlisting}

\noindent Executando na MEPA com entrada 10 via arquivo de entrada:
\begin{lstlisting}[language=bash]
python mepa/mepa_pt.py --progfile test_fib.mepa --infile input.txt --silent
\end{lstlisting}
\noindent Obteve-se como saída:
\begin{lstlisting}
205 instrucoes executadas

10
55
\end{lstlisting}
satisfazendo o critério oficial do Desafio 4.2 com absoluta exatidão.

\subsection{Desafio 4.3 -- Comparação Estruturada com o Compilador Java de Referência}

A Tabela~\ref{tab:correspondencia} detalha a correspondência formal entre as regras
implementadas em Python com PLY e o analisador sintático descendente recursivo
(\texttt{Parser.java}) do compilador $\mu$C de referência.

\begin{table}[H]
\centering
\caption{Mapeamento estruturado entre o compilador didático Java e o compilador PLY.}
\label{tab:correspondencia}
\small
\begin{tabular}{p{4.2cm}p{4.8cm}p{5.5cm}}
\toprule
\textbf{Construção Sintática} & \textbf{Método no \texttt{Parser.java}} & \textbf{Regra PLY e Instruções MEPA} \\
\midrule
Bloco mestre & \texttt{program()} & \texttt{p\_program} (\texttt{INPP}, \texttt{PARA}, \texttt{FIM}) \\
Declaração de tipos & \texttt{decls()} & \texttt{p\_decl\_int}, \texttt{p\_decl\_float} (\texttt{AMEM}) \\
Atribuição & \texttt{assign()} & \texttt{p\_stmt\_assign} (\texttt{ARMZ 0,n}) \\
Escrita em tela & \texttt{stmt()} (\texttt{Tag.PRINT}) & \texttt{p\_stmt\_print} (\texttt{IMPR}) \\
Leitura de teclado & \texttt{stmt()} (\texttt{Tag.READ}) & \texttt{p\_stmt\_read} (\texttt{LEIT}, \texttt{ARMZ}) \\
Condicional simples & \texttt{stmt()} (\texttt{Tag.IF}) & \texttt{p\_if\_head}, \texttt{p\_stmt\_if} (\texttt{DSVF}) \\
Condicional composto & \texttt{stmt()} (\texttt{Tag.ELSE}) & \texttt{p\_else\_mark}, \texttt{p\_stmt\_if\_else} (\texttt{DSVS}) \\
Laço pré-testado & \texttt{stmt()} (\texttt{Tag.WHILE}) & \texttt{p\_while\_start}, \texttt{p\_while\_head} (\texttt{DSVF}, \texttt{DSVS}) \\
Laço pós-testado & \texttt{stmt()} (\texttt{Tag.DO}) & \texttt{p\_do\_head}, \texttt{p\_stmt\_do} (\texttt{DSVF}, \texttt{DSVS}) \\
Interrupção & \texttt{stmt()} (\texttt{Tag.BREAK}) & \texttt{p\_stmt\_break} (\texttt{DSVS}) \\
Adição / Subtração & \texttt{expr()} & \texttt{p\_expr\_plus}, \texttt{p\_expr\_minus} (\texttt{SOMA}, \texttt{SUBT}) \\
Multiplicação / Divisão & \texttt{term()} & \texttt{p\_term\_times}, \texttt{p\_term\_divide} (\texttt{MULT}, \texttt{DIVI}) \\
Inversão aritmética & \texttt{unary()} & \texttt{p\_unary\_minus} (\texttt{INVR}) \\
Disjunção / Conjunção & \texttt{bool()}, \texttt{join()} & \texttt{p\_bool\_or}, \texttt{p\_bool\_and} (\texttt{DISJ}, \texttt{CONJ}) \\
Comparações relacionais & \texttt{rel()} & \texttt{p\_relation\_cmp} (\texttt{CMIG}, \texttt{CMDG}, \texttt{CMME}, \dots) \\
\bottomrule
\end{tabular}
\end{table}

\subsection{Desafio 4.4 -- Auditoria Crítica de Código Gerado por IA}

Na Seção 12.4 do roteiro, é submetido à auditoria o seguinte trecho de código
proposto por um modelo generativo de linguagem:

\begin{lstlisting}[language=Python]
def p_stmt_while(p):
    "stmt : WHILE '(' bool ')' stmt"
    l1 = new_label()
    emit(f"L{l1}:")        # (A)
    emit(f"DSVF L{l2}")    # (B)
    emit(f"DSVS L{l1}")    # (C)
    emit(f"L{l2}:")        # (D)
\end{lstlisting}

\subsubsection{Identificação dos Erros Críticos}

A auditoria identificou \textbf{dois erros conceituais e de execução}:

\begin{enumerate}
    \item \textbf{Erro 1 -- Falha de Escopo e Variável Inexistente (\texttt{NameError}):}
          A variável \texttt{l2} é utilizada nas linhas (B) e (D) sem nunca ter sido
          declarada ou instanciada. O código gerado executa apenas \texttt{l1 = new\_label()}.
          Em Python, isso provoca uma interrupção fatal de execução (\texttt{NameError: name 'l2' is not defined}).
    \item \textbf{Erro 2 -- Inversão Temporal de Emissão (Ordem Incorreta das Reduções):}
          Em analisadores LALR, a função associada a uma regra só executa quando todos os
          símbolos da produção já foram reduzidos. Isso implica que todo o código MEPA
          de avaliação da condição (\texttt{bool}) e todo o código do corpo do laço
          (\texttt{stmt}) já foram emitidos no fluxo de saída!
          Assim, emitir \texttt{L1:}, \texttt{DSVF L2}, \texttt{DSVS L1} e \texttt{L2:}
          dentro de uma única função pós-redução faz com que todas essas instruções
          sejam agrupadas em bloco \textbf{depois} do corpo do laço. O rótulo inicial
          fica posicionado após o corpo, o desvio condicional não protege a execução,
          e o laço entra em comportamento caótico.
\end{enumerate}

\subsubsection{Por que o Código Parecia Plausível?}

O código gerado pela IA possui alta plausibilidade superficial porque ele
reproduz fielmente o padrão visual e descritivo encontrado em livros-texto de
compiladores (como o clássico ``Livro do Dragão'' de Aho et al.):
\begin{center}
\textit{rótulo\_início} $\rightarrow$ \textit{desvio\_falso} $\rightarrow$ \textit{desvio\_retorno} $\rightarrow$ \textit{rótulo\_fim}
\end{center}
A falha do modelo reside na incapacidade de relacionar a \textbf{sintaxe textual}
com a \textbf{dinâmica temporal} de um autômato ascendente com ações pós-redução.

\subsubsection{Versão Corrigida e Validação}

A versão correta exige desmembrar a regra em ações prévias (\texttt{while\_start} e
\texttt{while\_head}), conforme demonstrado na Seção 7 deste relatório e comprovado
por execução bem-sucedida na MEPA.

% =========================================================
\section{Comparação Detalhada com o Compilador de Referência}
% =========================================================

\subsection{Matriz Comparativa das Sequências de Instruções}

Compilaram-se os mesmos programas-fonte em ambos os sistemas. A Tabela~\ref{tab:comparacao}
apresenta o confronto direto das instruções emitidas.

\begin{longtable}{p{2.6cm}p{6.4cm}p{6.4cm}}
\caption{Comparação de código MEPA: Compilador Desenvolvido vs Compilador Java de Referência.}
\label{tab:comparacao}\\
\toprule
\textbf{Programa} & \textbf{Compilador Desenvolvido (PLY)} & \textbf{Compilador de Referência ($\mu$C Java)} \\
\midrule
\endfirsthead
\toprule
\textbf{Programa} & \textbf{Compilador Desenvolvido (PLY)} & \textbf{Compilador de Referência ($\mu$C Java)} \\
\midrule
\endhead
\bottomrule
\endfoot

\textbf{Expressão} &
\texttt{INPP}\newline
\texttt{AMEM 4 ; a}\newline
\texttt{CRCT 1}\newline
\texttt{CRCT 2}\newline
\texttt{SOMA}\newline
\texttt{ARMZ 0,0 ; a}\newline
\texttt{PARA}\newline
\texttt{FIM} &
\texttt{INPP}\newline
\texttt{AMEM 4 ; a}\newline
\texttt{AMEM 4 ; temp (ponteiro)}\newline
\texttt{CRCT 1}\newline
\texttt{CRCT 2}\newline
\texttt{SOMA}\newline
\texttt{ARMZ 0,0 ; a}\newline
\texttt{DMEM 8}\newline
\texttt{PARA}\newline
\texttt{FIM} \\
\midrule

\textbf{Condicional \texttt{if-else}} &
\texttt{INPP}\newline
\texttt{AMEM 4 ; a}\newline
\texttt{CRCT 5}\newline
\texttt{ARMZ 0,0}\newline
\texttt{CRVL 0,0}\newline
\texttt{CRCT 0}\newline
\texttt{CMMA}\newline
\texttt{DSVF L1}\newline
\texttt{CRCT 1}\newline
\texttt{ARMZ 0,0}\newline
\texttt{DSVS L2}\newline
\texttt{L1: NADA}\newline
\texttt{CRCT 2}\newline
\texttt{ARMZ 0,0}\newline
\texttt{L2: NADA}\newline
\texttt{CRVL 0,0}\newline
\texttt{IMPR}\newline
\texttt{PARA}\newline
\texttt{FIM} &
\texttt{INPP}\newline
\texttt{AMEM 4 ; a}\newline
\texttt{AMEM 4 ; temp (ponteiro)}\newline
\texttt{CRCT 5}\newline
\texttt{ARMZ 0,0}\newline
\texttt{CRVL 0,0}\newline
\texttt{CRCT 0}\newline
\texttt{CMMA}\newline
\texttt{DSVF L1}\newline
\texttt{CRCT 1}\newline
\texttt{ARMZ 0,0}\newline
\texttt{DSVS L2}\newline
\texttt{L1:}\newline
\texttt{CRCT 2}\newline
\texttt{ARMZ 0,0}\newline
\texttt{L2:}\newline
\texttt{CRVL 0,0}\newline
\texttt{IMPR}\newline
\texttt{DMEM 8}\newline
\texttt{PARA}\newline
\texttt{FIM} \\
\midrule

\textbf{Laço \texttt{while}} &
\texttt{INPP}\newline
\texttt{AMEM 4 ; k}\newline
\texttt{CRCT 1}\newline
\texttt{ARMZ 0,0}\newline
\texttt{L1: NADA}\newline
\texttt{CRVL 0,0}\newline
\texttt{CRCT 5}\newline
\texttt{CMEG}\newline
\texttt{DSVF L2}\newline
\texttt{CRVL 0,0}\newline
\texttt{IMPR}\newline
\texttt{CRVL 0,0}\newline
\texttt{CRCT 1}\newline
\texttt{SOMA}\newline
\texttt{ARMZ 0,0}\newline
\texttt{DSVS L1}\newline
\texttt{L2: NADA}\newline
\texttt{PARA}\newline
\texttt{FIM} &
\texttt{INPP}\newline
\texttt{AMEM 4 ; k}\newline
\texttt{AMEM 4 ; temp (ponteiro)}\newline
\texttt{CRCT 1}\newline
\texttt{ARMZ 0,0}\newline
\texttt{L1:}\newline
\texttt{CRVL 0,0}\newline
\texttt{CRCT 5}\newline
\texttt{CMEG}\newline
\texttt{DSVF L2}\newline
\texttt{AMEM 4 ; temp (ponteiro)}\newline
\texttt{CRVL 0,0}\newline
\texttt{IMPR}\newline
\texttt{CRVL 0,0}\newline
\texttt{CRCT 1}\newline
\texttt{SOMA}\newline
\texttt{ARMZ 0,0}\newline
\texttt{DMEM 4}\newline
\texttt{DSVS L1}\newline
\texttt{L2:}\newline
\texttt{DMEM 8}\newline
\texttt{PARA}\newline
\texttt{FIM} \\
\midrule

\textbf{Interrupção \texttt{break}} &
\texttt{...}\newline
\texttt{CRVL 0,0}\newline
\texttt{CRCT 5}\newline
\texttt{CMMA}\newline
\texttt{DSVF L3}\newline
\texttt{\textbf{DSVS L2 ; break}}\newline
\texttt{L3: NADA}\newline
\texttt{CRVL 0,0}\newline
\texttt{IMPR}\newline
\texttt{...} &
\texttt{...}\newline
\texttt{CRVL 0,0}\newline
\texttt{CRCT 5}\newline
\texttt{CMMA}\newline
\texttt{DSVF L3}\newline
\textit{[DSVS omitido por falha]}\newline
\texttt{L3:}\newline
\texttt{CRVL 0,0}\newline
\texttt{IMPR}\newline
\texttt{...} \\
\end{longtable}

\subsection{Análise das Divergências Encontradas}

A análise comparativa evidencia quatro divergências principais:

\begin{enumerate}
    \item \textbf{Slot de Ponteiro Temporário (\texttt{AMEM 4 ; temp}):} O gerador Java
          reserva compulsoriamente um slot temporário de 4 bytes a cada bloco para dar
          suporte a atribuições indiretas via arranjos e ponteiros (\texttt{Env.reserveTemp}).
          Como o nosso compilador implementou a especificação de tradução do laboratório,
          ele emite apenas o espaço requerido pelas variáveis declaradas.
    \item \textbf{Desalocação Dinâmica com \texttt{DMEM}:} O compilador Java emite
          instruções \texttt{DMEM} ao final de cada bloco de comandos para desempilhar
          variáveis locais. O nosso compilador adota o modelo estático de memória plana
          com alocação consolidada no prólogo, garantindo segurança contra estouro de
          pilha em laços.
    \item \textbf{Rótulos em Linha Própria vs \texttt{label: NADA}:} O compilador Java
          escreve rótulos isolados, demandando pós-processamento externo para a MEPA.
          O nosso compilador emite \texttt{L1: NADA}, sendo imediatamente executável.
    \item \textbf{Correção Semântica do \texttt{break}:} O compilador de referência
          omite o desvio \texttt{DSVS} do break. O nosso compilador gera a instrução
          correta, garantindo a saída imediata do laço.
\end{enumerate}

% =========================================================
\section{Decisões de Implementação}
% =========================================================

\subsection{Dicionário de Palavras Reservadas vs Expressões Regulares Específicas}

Decidiu-se mapear palavras-chave através do dicionário \texttt{reserved} em conjunto
com a regra genérica de identificadores \texttt{t\_ID}.

\noindent\textit{Alternativa descartada:} Criar regras específicas como \texttt{t\_IF = r'if'}.

\noindent\textit{Justificativa:} A criação de regras dedicadas aumenta desnecessariamente
o autômato léxico e cria armadilhas com identificadores que possuem palavras reservadas
como prefixo (como \texttt{ifx} ou \texttt{while1}). O uso do dicionário garante tempo
$O(1)$ de consulta e total previsibilidade.

\subsection{Alocação no Prólogo vs Emissão no Ponto da Declaração}

Optou-se por acumular as instruções de alocação de memória na lista \texttt{ALLOCATIONS}
e emiti-las todas no prólogo, imediatamente após \texttt{INPP}.

\noindent\textit{Alternativa descartada:} Emitir \texttt{AMEM} no corpo da produção
no exato instante em que a declaração é reduzida.

\noindent\textit{Justificativa:} Caso uma declaração ocorra dentro de um bloco aninhado
que esteja no interior de um \texttt{while} ou \texttt{do-while}, emitir \texttt{AMEM}
no ponto da redução repetiria a alocação a cada iteração do laço, estourando a pilha da MEPA.

\subsection{Tabela de Símbolos com Pilha de Escopos}

Adotou-se uma pilha de dicionários (\texttt{SCOPES = []}) para representar o escopo
léxico, complementada pelo controle linear de \texttt{OFFSET}.

\noindent\textit{Alternativa descartada:} Dicionário simples sem controle de escopo.

\noindent\textit{Justificativa:} A pilha de escopos permite sombreamento legítimo de variáveis
em blocos aninhados, enquanto rejeita re-declarações ilegais dentro do mesmo bloco e
rejeita acessos a variáveis fora de escopo.

\subsection{Emissão Integrada com \texttt{NADA} para MEPA}

Adotou-se o formato \texttt{mark\_label(label)} emitindo \texttt{f"{label}: NADA"}.

\noindent\textit{Alternativa descartada:} Depender do script \texttt{awk} sugerido na Seção 11 do guia.

\noindent\textit{Justificativa:} A eliminação de dependências externas torna o pipeline de
compilação autosuficiente e portátil entre diferentes sistemas operacionais (Windows/Linux).

% =========================================================
\section{Evidências Experimentais e Comandos de Reprodução}
% =========================================================

Todos os testes podem ser reproduzidos na raiz do projeto com os seguintes comandos:

\subsection{Teste 1: Expressões e Precedência}
\begin{lstlisting}[language=bash]
python lab1_tds.py exemplos/expr.uc -o test_expr.mepa
python mepa/mepa_pt.py --progfile test_expr.mepa --silent
# Resultado: Executa 7 instrucoes com sucesso.
\end{lstlisting}

\subsection{Teste 2: Condicional \texttt{if-else}}
\begin{lstlisting}[language=bash]
python lab1_tds.py exemplos/if_else.uc -o test_if.mepa
python mepa/mepa_pt.py --progfile test_if.mepa --silent
# Saida observada: 1 (para a = 5) e 2 (para a = -5).
\end{lstlisting}

\subsection{Teste 3: Laço \texttt{while}}
\begin{lstlisting}[language=bash]
python lab1_tds.py exemplos/while.uc -o test_while.mepa
python mepa/mepa_pt.py --progfile test_while.mepa --silent
# Saida observada: 1, 2, 3, 4, 5.
\end{lstlisting}

\subsection{Teste 4: Interrupção por \texttt{break}}
\begin{lstlisting}[language=bash]
python lab1_tds.py exemplos/break.uc -o test_break.mepa
python mepa/mepa_pt.py --progfile test_break.mepa --silent
# Saida observada: 1, 2, 3, 4, 5 (laco interrompido ao atingir k > 5).
\end{lstlisting}

\subsection{Teste 5: Desafio 4.1 (\texttt{do-while})}
\begin{lstlisting}[language=bash]
python lab1_tds.py exemplos/do_while.uc -o test_do_while.mepa
python mepa/mepa_pt.py --progfile test_do_while.mepa --silent
# Saida observada: 1, 2, 3.
\end{lstlisting}

\subsection{Teste 6: Desafio 4.2 (Fibonacci com \texttt{read} e Entrada 10)}
\begin{lstlisting}[language=bash]
python lab1_tds.py exemplos/fib.uc -o test_fib.mepa
python mepa/mepa_pt.py --progfile test_fib.mepa --infile input.txt --silent
# Saida observada: 10 e 55 (205 instrucoes MEPA).
\end{lstlisting}

% =========================================================
\section{Dificuldades Encontradas e Soluções}
% =========================================================

\begin{enumerate}
    \item \textbf{Atributos Herdados no Modelo Ascendente:} Compreender por que o desvio
          condicional \texttt{DSVF} ficava posicionado após o corpo do \texttt{if} quando
          emitido na redução da regra foi a principal dificuldade teórica do trabalho.
          A superação dessa barreira exigiu internalizar o conceito de ações semânticas
          intermediárias (\textit{mid-rule actions}) simuladas por regras vazias e
          pilhas de controle globais (\texttt{IF\_PENDING}, \texttt{WHILE\_PENDING}).
    \item \textbf{Alocação de Memória dentro de Laços:} Ao testar blocos aninhados contendo
          declarações dentro de repetições, notou-se o perigo da multiplicação descontrolada
          de instruções \texttt{AMEM}. A solução concebida pela equipe (prolongar a emissão
          de \texttt{AMEM} para o prólogo) garantiu estabilidade e conformidade total
          com a arquitetura MEPA.
    \item \textbf{Particularidade da MEPA com Ponto Flutuante:} Descobriu-se durante os
          testes que o carregador da MEPA 5.0 (\texttt{fixArgs} em \texttt{mepa\_defs.py})
          valida os argumentos das instruções utilizando o método \texttt{a.isdigit()},
          rejeitando constantes fracionárias literais como \texttt{0.0314}. O compilador
          calcula e aloca perfeitamente os 8 bytes de \texttt{float}, mas para simulações
          completas na MEPA os cálculos devem utilizar dados inteiros.
\end{enumerate}

% =========================================================
\section{Conclusão}
% =========================================================

A realização do Laboratório 1 permitiu consolidar de forma prática e aprofundada os
fundamentos de compilação estudados ao longo dos Módulos 1, 2, 3 e 4. A utilização
do gerador PLY demonstrou a viabilidade de construir um compilador funcional e
robusto com código enxuto e de alta legibilidade.

A técnica de Tradução Dirigida por Sintaxe revelou sua elegância ao permitir que a
própria estrutura da derivação sintática guie a emissão do código executável da
máquina de pilha. O domínio dos conceitos de atributos sintetizados e herdados foi
determinante para solucionar os desafios de controle de fluxo (\texttt{if}, \texttt{while},
\texttt{do-while} e \texttt{break}).

Por fim, a comparação minuciosa contra o compilador Java de referência e a auditoria
crítica de código gerado por IA comprovaram que o entendimento teórico dos mecanismos
internos de um compilador é absolutamente insubstituível para projetar, depurar e
garantir a segurança de linguagens e ferramentas de software.

% =========================================================
\section{Declaração de Uso de Inteligência Artificial (Escala AIAS)}
% =========================================================

Em estrita consonância com a Política Institucional de Uso de Inteligência Artificial
da UNIFEI e com a Seção 14 do guia da disciplina, os integrantes da equipe declaram:

\subsection*{Aluno 1}
Declaro que utilizei ferramentas de inteligência artificial estritamente dentro dos
limites autorizados pela política da disciplina (Escala AIAS). Nos Passos 1 a 3, a IA
foi consultada em nível auxiliar (\textbf{AIAS Nível 2}) para esclarecimento de mensagens
de erro do PLY e verificação de regras de formatação. Os núcleos conceituais dos
Passos 4 e 5 (\textit{mid-rule actions}, atributos herdados e pilhas de desvio) foram
concebidos, codificados e validados de forma integralmente autoral pela equipe
(\textbf{AIAS Nível 1 -- Sem IA}). No Desafio 4.4, atuei como auditor crítico
(\textbf{AIAS Nível 4}), diagnosticando as falhas da solução proposta pelo modelo e
fundamentando o veredito técnico com base na teoria de autômatos ascendentes.

\subsection*{Aluno 2}
Declaro que respeitei integralmente a política de uso de IA estipulada para a atividade.
O desenvolvimento dos mecanismos de controle de fluxo (condicionais e laços) e o
desenho das estruturas de dados de escopo foram resolvidos de forma estritamente autoral
com os colegas de grupo (\textbf{AIAS Nível 1}), constituindo evidência concreta de
aprendizado. As consultas a ferramentas generativas restringiram-se a aspectos de
sintaxe do PLY e revisão tipográfica do texto em \LaTeX{} (\textbf{AIAS Nível 2}).

% =========================================================
\section{Referências Bibliográficas}
% =========================================================

\begin{thebibliography}{9}

\bibitem{aho2008}
AHO, Alfred V.; LAM, Monica S.; SETHI, Ravi; ULLMAN, Jeffrey D.
\textit{Compiladores: princípios, técnicas e ferramentas}.
2. ed. São Paulo: Pearson, 2008.

\bibitem{ply}
BEASLEY, David.
\textit{PLY (Python Lex-Yacc) Documentation}.
Disponível em: \url{https://www.dabeaz.com/ply/}. Acesso em: out. 2026.

\bibitem{louden}
LOUDEN, Kenneth C.
\textit{Compiladores: princípios e práticas}.
São Paulo: Pioneira Thomson Learning, 2004.

\bibitem{mepa}
KOWALTOWSKI, Tomasz.
\textit{MEPA 5.0 -- Interpretador e especificação das instruções}.
Instituto de Computação, Universidade Estadual de Campinas (UNICAMP), 2015.

\end{thebibliography}

\end{document}